\begingroup%
\def\baselinestretch{1}%
\definecolor{D15Teal}{HTML}{27877F}%
\definecolor{D15Amber}{HTML}{B7791F}%
\definecolor{D15Ink}{HTML}{243447}%
\definecolor{D15Slate}{HTML}{64748B}%
\definecolor{D15PaleTeal}{HTML}{EFF8F5}%
\def\Gap#1{\textcolor{D15Amber}{#1}}%
\begin{tikzpicture}[
  x=1mm, y=1mm,
  rung/.style={draw=D15Slate!60, fill=white, rounded corners=1.8mm,
    line width=.6pt, minimum width=152mm, minimum height=10.5mm, outer sep=0pt},
  name/.style={text=D15Ink, align=center, text width=22mm,
    font=\bfseries\fontsize{9.5}{11.5}\selectfont},
  cell/.style={text=D15Ink, align=center, font=\fontsize{8.5}{10.5}\selectfont},
  head/.style={text=D15Ink, align=center,
    font=\bfseries\fontsize{9}{11}\selectfont},
  note/.style={text=D15Slate, font=\fontsize{8.5}{10.5}\selectfont},
  flow/.style={-{Stealth[length=2.25mm,width=1.65mm]}, line width=.7pt}
]
  \def\cA{12}\def\cB{48}\def\cC{93}\def\cD{134}%
  \node[head] at (\cA,61) {层级};
  \node[head] at (\cB,61) {设计层\\{\mdseries（Schema 的对应物）}};
  \node[head] at (\cC,61) {验证器\\{\mdseries（单次反馈周期）}};
  \node[head] at (\cD,61) {成本先验\\{\mdseries（Microbench 的对应物）}};

  \node[rung, draw=D15Teal, fill=D15PaleTeal, line width=.9pt] at (76,0) {};
  \node[name] at (\cA,0) {Kernel\\[-.3mm]{\mdseries\fontsize{7.5}{9}\selectfont\color{D15Teal}本文起点}};
  \node[cell, text width=44mm] at (\cB,0) {数学改写 · 分块映射 · 流水编排};
  \node[cell, text width=46mm] at (\cC,0) {编译 + 数值对齐 + 实测延迟\\（秒～分钟）};
  \node[cell, text width=34mm] at (\cD,0) {Microbench\\Roofline};

  \node[rung] at (76,12) {};
  \node[name] at (\cA,12) {Megakernel};
  \node[cell, text width=44mm] at (\cB,12) {融合边界 · 阶段交接\\跨算子调度};
  \node[cell, text width=46mm] at (\cC,12) {端到端数值对齐 + 单步延迟\\（分钟）};
  \node[cell, text width=34mm] at (\cD,12) {Microbench\\\Gap{+ 跨算子干扰（待建）}};

  \node[rung] at (76,24) {};
  \node[name] at (\cA,24) {推理框架};
  \node[cell, text width=44mm] at (\cB,24) {并行策略 · 批调度\\PD 分离 · KV Cache 管理};
  \node[cell, text width=46mm] at (\cC,24) {与参考输出一致 + 吞吐 / SLO\\（分钟～小时）};
  \node[cell, text width=34mm] at (\cD,24) {计算与通信 Roofline\\\Gap{真实负载模型（待建）}};

  \node[rung] at (76,36) {};
  \node[name] at (\cA,36) {训练框架};
  \node[cell, text width=44mm] at (\cB,36) {并行切分 · 重计算\\通信重叠 · 精度选择};
  \node[cell, text width=46mm] at (\cC,36) {数值与 Loss 曲线对齐 + MFU\\（小时～天）};
  \node[cell, text width=34mm] at (\cD,36) {\Gap{保序的小规模}\\\Gap{代理实验（待建）}};

  \node[rung] at (76,48) {};
  \node[name] at (\cA,48) {自动化科研};
  \node[cell, text width=44mm] at (\cB,48) {研究假设 · 实验设计};
  \node[cell, text width=46mm] at (\cC,48) {\Gap{可复现的客观评测（多数尚缺）}\\（天～月）};
  \node[cell, text width=34mm] at (\cD,48) {\Gap{尚缺}};

  \draw[flow, D15Amber] (-5,-5) -- (-5,53);
  \node[note, text=D15Ink, rotate=90] at (-9,24) {越往上，设计空间越大，反馈越慢};

  \node[note, anchor=west] at (0,-9.5) {\Gap{琥珀色}：迁移到该层级时尚需建立的部分};
  \node[note, anchor=east] at (152,-9.5) {实现层分别对应 CUDA / CuTe 代码、框架代码与配置、实验代码};
\end{tikzpicture}%
\endgroup%
